\originalsection{线性三原子分子的伸缩与弯曲}
\begin{example}
\textbf{CO$_2$的弹性模型.} 
两个端点原子质量均为$m$, 中间原子质量为$M$, 平衡位置沿$x$轴依次为$-b,0,b$. 两键的伸缩弹性常数为$k>0$, 每一横向弯曲方向的有效常数为$\kappa>0$. 求小振动频率与质量归一化模, 并解释振动数目.
\end{example}
\begin{solution}
先考虑纵向位移$x_1,x_2,x_3$, 它们表示相对平衡位置的偏移, 不是原子绝对坐标. 到二阶,
\[
T_x=\tfrac12(m\dot x_1^2+M\dot x_2^2+m\dot x_3^2),\qquad
U_x=\tfrac{k}{2}(x_2-x_1)^2+\tfrac{k}{2}(x_3-x_2)^2.
\]
所以质量矩阵$\mathcal M=\operatorname{diag}(m,M,m)$, 刚度矩阵为
\[
\mathcal K_x=k\begin{pmatrix}1&-1&0\\-1&2&-1\\0&-1&1\end{pmatrix}.
\]
令$\rho=m/M$. 整体平移向量$(1,1,1)$是零模, 而两个伸缩向量可由端点交换对称性构造. 对$(1,0,-1)$直接乘矩阵得到$\mathcal K_x(1,0,-1)^{\mathsf T}=k(1,0,-1)^{\mathsf T}$, 中间原子不动, 因而$\omega_-^2=k/m$. 对$(1,-2\rho,1)$, 质心偏移$m-2M\rho+m=0$, 乘矩阵为$k(1+2\rho)(1,-2,1)^{\mathsf T}$; 同时质量矩阵乘该向量为$m(1,-2,1)^{\mathsf T}$. 因而$\omega_+^2=k(1+2\rho)/m$.

在质量内积$\langle v,w\rangle_{\mathcal M}=v^{\mathsf T}\mathcal M w$下, 三个归一化向量为
\[
a_0=\frac{(1,1,1)^{\mathsf T}}{\sqrt{2m+M}},\quad
 a_- =\frac{(1,0,-1)^{\mathsf T}}{\sqrt{2m}},\quad
 a_+ =\frac{(1,-2\rho,1)^{\mathsf T}}{\sqrt{2m(1+2\rho)}}.
\]
它们两两质量正交. 写$x=a_0Q_0+a_-Q_-+a_+Q_+$后, 动能为$\sum\dot Q_j^2/2$, 势能为$(\omega_-^2Q_-^2+\omega_+^2Q_+^2)/2$. $Q_0$是质心平移, 应与内部振动分开; 初值投影为$Q_j=a_j^{\mathsf T}\mathcal M x$. 归一化使$Q_j$具有$\sqrt{\text{质量}}\times\text{长度}$的单位, 不应叫无量纲坐标.

只有键长弹簧时, 横向偏移造成的键长变化是二阶小量, 因而伸缩能的横向回复项是四阶, 不会给有限线性弯曲频率. 必须另加角弯曲势. 在选定$xy$平面内, 定义横向位移$y_i$, 两键倾角到一阶分别为$(y_2-y_1)/b$与$(y_3-y_2)/b$. 两倾角之差与
$D_y=y_1-2y_2+y_3$成正比. 取旋转不变的二次弯曲模型
\[
U_y=\frac{\kappa}{4}(y_1-2y_2+y_3)^2,\qquad
\mathcal K_y=\frac{\kappa}{2}\vect c\vect c^{\mathsf T},\quad
\vect c=(1,-2,1)^{\mathsf T}.
\]
这里$\kappa$单位同$k$; 若直接以两键夹角偏离的平方写$U=C(\Delta\alpha)^2/2$, 则$\kappa=2C/b^2$, 不应遗漏长度尺度. 平移$(1,1,1)$和刚体小转动$(-b,0,b)$都满足$D_y=0$, 因而是两个零模. 唯一平面内弯曲模为
\[
a_b=\frac{(1,-2\rho,1)^{\mathsf T}}{\sqrt{2m(1+2\rho)}},\qquad
\omega_b^2=\frac{\kappa}{m}(1+2\rho).
\]
核验时$\vect c^{\mathsf T}(1,-2\rho,1)=2(1+2\rho)$, 所以$\mathcal K_y(1,-2\rho,1)^{\mathsf T}=\kappa(1+2\rho)\vect c$, 与$\omega_b^2\mathcal M(1,-2\rho,1)^{\mathsf T}$一致. 这个$a_b$与$a_+$虽然有相同的原子振幅比例, 作用的空间方向不同, 是不同的物理模.

与8.09第57--59页的平面约化对应时, 固定质心并取平面姿态后可令$y_1=y_3=y$, $y_2=-2\rho y$, 纵向端点偏移为$\eta_1,\eta_3$, 中间偏移为$-\rho(\eta_1+\eta_3)$. 用$\eta=(\eta_1,y,\eta_3)$写,
\begin{align*}
M_{\rm red}&=m\begin{pmatrix}1+\rho&0&\rho\\0&2(1+2\rho)&0\\\rho&0&1+\rho\end{pmatrix},\\
K_{\rm red}&=\begin{pmatrix}k(1+2\rho+2\rho^2)&0&2k\rho(1+\rho)\\0&2\kappa(1+2\rho)^2&0\\2k\rho(1+\rho)&0&k(1+2\rho+2\rho^2)\end{pmatrix}.
\end{align*}
两纵向模分别为$(1,0,-1)/\sqrt{2m}$与$(1,0,1)/\sqrt{2m(1+2\rho)}$, 横向模为$(0,1,0)/\sqrt{2m(1+2\rho)}$. 这给出与全坐标法相同的三个平面频率, 也解释了原稿的质量矩阵为什么不是单位阵.
\end{solution}

\begin{figure}[htbp]
\centering
\begin{tikzpicture}[>=stealth,scale=.9]
\foreach \y/\lab in {0/伸缩一,1.5/伸缩二,3/平面弯曲} {
\draw (-2,\y)--(0,\y)--(2,\y);
\fill (-2,\y) circle (3pt);\fill (0,\y) circle (4pt);\fill (2,\y) circle (3pt);
\node[left] at (-3,\y) {\lab};}
\draw[->,thick] (-2,0)--(-1.3,0);\draw[->,thick] (2,0)--(1.3,0);
\draw[->,thick] (-2,1.5)--(-1.3,1.5);\draw[->,thick] (2,1.5)--(2.7,1.5);\draw[->,thick] (0,1.5)--(-1,1.5);
\draw[->,thick] (-2,3)--(-2,3.65);\draw[->,thick] (2,3)--(2,3.65);\draw[->,thick] (0,3)--(0,2.35);
\node at (-2,-.55) {$m$};\node at (0,-.55) {$M$};\node at (2,-.55) {$m$};
\end{tikzpicture}
\caption{三个平面内部模的方向示意, 箭头长度不表示质量归一化振幅. 第二伸缩及弯曲的中间原子振幅为端点的$-2m/M$倍; 弯曲还有一个垂直纸面的简并伙伴.}
\end{figure}

三维线性分子共有9个位移坐标. 整体平移有3个独立方向, 但绕分子轴的小转动不改变平衡点的位置, 只有垂直于分子轴的两个转动能产生独立一阶位移. 因而零模数为5, 内部振动为$9-5=4$, 即两个伸缩及两个简并弯曲. 在上面的模型中再加$U_z=\kappa(z_1-2z_2+z_3)^2/4$, 就得到第二弯曲模与同样的$\omega_b$. 两个横向方向构成二维简并子空间, 可以任取正交基; 圆偏振组合也只是同一子空间中两个实模的相位组合.

原8.09第56页写“9个坐标减6个刚体坐标剩3个振动”, 对三维线性平衡构型不适用. 其后选定平面的约化三模计算可保留, 但不能当作三维全部振动. 该弹性模型描述模态结构, 并未用真实分子光谱拟合$k,\kappa$, 因此这些符号频率不能直接冒充实验CO$_2$波数.

新增内容来源: 8.09第3章第56--59页; 本册以全坐标质量内积重新推导模态并纠正三维零模计数, 弯曲势采用与其平面二次型一致且保留刚体转动零模的形式.
